\scanpage{148}
% archive foliation: 147
vu que
\[
g^{-1}=\begin{pmatrix}d&-b\\-c&a\end{pmatrix}
\]
la formule
\[
\tau\circ g^{-1}=\frac{a\tau-b}{-c\tau+d}
\tag{11}
\]
Introduisons l'élément
\[
\left\{
\begin{aligned}
&\tau_\infty=\begin{pmatrix}-1&0\\0&+1\end{pmatrix}\in\mathrm{Gl}(2,\mathbb Z)
&&\text{d'où}\\
&\tau_\infty\begin{pmatrix}A&B\\C&D\end{pmatrix}\tau_\infty^{-1}
=\tau_\infty\left(\begin{pmatrix}A&B\\C&D\end{pmatrix}\right)
=\begin{pmatrix}A&-B\\-C&D\end{pmatrix}\\
&\tau_\infty\left(\begin{pmatrix}a&b\\c&d\end{pmatrix}^{-1}\right)
=\begin{pmatrix}d&b\\c&a\end{pmatrix}
&&\text{``échange de $a$ et de $d$''},
\end{aligned}
\right.
\tag{12}
\]
on trouve
\[
\tau\circ\tau_\infty(g^{-1})=\frac{a\tau+b}{c\tau+d}
\qquad
\begin{gathered}
\text{où }g=\begin{pmatrix}a&b\\c&d\end{pmatrix}\\
\tau_\infty(g^{-1})=\begin{pmatrix}d&b\\c&a\end{pmatrix}
\end{gathered}
\tag{13}
\]
Par la suite, quand on fait opérer $\mathrm{Sl}(2,\mathbb Z)$
sur $\mathcal D^+$, ce sera tjrs (sauf mention expresse)
par la formule (13) --- ce n'est donc pas
l'action ``modulaire'' la plus évidente, mais
celle consacrée par l'usage, qui correspond
à la précédente \inserted{(à droite)}, modifiée via le
passage
\[
\underset{\begin{pmatrix}a&b\\c&d\end{pmatrix}}{g}
\longmapsto\tau_\infty(g^{-1})=\begin{pmatrix}d&b\\c&a\end{pmatrix}.
\]
